% GENERATED FILE. Edit canonical lessons, then run npm run generate:books.
% canonical-source-hash: fnv1a64:102ec5aa740022ee
% canonical-lessons: BN-C479-ghum, BN-C479-onidra, BN-C479-tondra, BN-C479-jhimano, BN-C479-armora-bhanga

\chapter{Sleep}
\label{ch:bn-sleep}
\hlchaptermodality{479}

\begin{chapteropening}
\textbf{By the end of this chapter:} \emph{I can talk about sleep, sleeplessness and dozing off.}

Complete the last lesson of chapter 479: I can talk about sleep, sleeplessness and dozing off.
\end{chapteropening}

\section[\texorpdfstring{\bn{ঘুম} (ghum)}{ঘুম (ghum)}]{\bn{ঘুম} (ghum) — sleep}

\label{lesson:BN-C479-ghum}



\begin{quote}
Before the new one: say the Bengali for a cold compress, then the Bengali for to nurse.
\end{quote}

\subsection*{You'll want to know: \bn{ঘুম}}
\textbf{\bn{ঘুম}} — \emph{ghum} — \textquotedblleft{}sleep\textquotedblright{}.

\begin{quote}
\bn{কাল} \bn{রাতে} \bn{আমার} \bn{ভালো} \bn{ঘুম} \bn{হল।}
\end{quote}

\begin{quote}
\emph{kāl rāte āmār bhālo ghum hôlô.}
\end{quote}

\begin{quote}
I slept well last night.
\end{quote}

\bn{ঘুম} \bn{পাওয়া}, to feel sleepy: \bn{আমার} \bn{ঘুম} \bn{পাচ্ছে}.

\begin{cousinweb}
\textbf{\bn{ঘুম}} (\emph{ghum}) is listed by CDIAL for Bengali, as dozing, under a reconstructed verb of turning round (CDIAL 4485); it shows the new letter \bn{ঘ}.
\end{cousinweb}

\begin{grammarlens}[title={What you've built}]
One more word for describing what you lived and felt.
\end{grammarlens}

\subsection*{Your turn}
\begin{itemize}
\raggedright
  \item \emph{Say it:} \emph{ghum}
  \item \emph{Say it:} \emph{ghum}, once more
  \item \emph{Say it:} \emph{kāl rāte āmār bhālo ghum hôlô.} — I slept well last night.
  \item \emph{Recall:} say the Bengali for a cold compress, then the Bengali for to nurse, then say \emph{ghum} again
\end{itemize}

\begin{tcolorbox}[breakable,colback=teal!4,colframe=teal!35!black,title={Before you move on}]
What does \bn{ঘুম} mean? (\textquotedblleft{}Sleep\textquotedblright{}.) Say it once more.
\end{tcolorbox}
\section[\texorpdfstring{\bn{অনিদ্রা} (ônidrā)}{অনিদ্রা (ônidrā)}]{\bn{অনিদ্রা} (ônidrā) — sleeplessness; insomnia}

\label{lesson:BN-C479-onidra}



\begin{quote}
Before the new one: say the Bengali for to nurse, then the Bengali for sleep.
\end{quote}

\subsection*{You'll want to know: \bn{অনিদ্রা}}
\textbf{\bn{অনিদ্রা}} — \emph{ônidrā} — \textquotedblleft{}sleeplessness; insomnia\textquotedblright{}.

\begin{quote}
\bn{চিন্তায়} \bn{বাবা} \bn{অনিদ্রায়} \bn{ভুগছেন।}
\end{quote}

\begin{quote}
\emph{chintāy bābā ônidrāy bhugchhen.}
\end{quote}

\begin{quote}
Father is suffering from insomnia because of worry.
\end{quote}

\begin{cousinweb}
\textbf{\bn{অনিদ্রা}} (\emph{ônidrā}) is a tatsama, Sanskrit \emph{anidrā}, \textquotedblleft{}sleeplessness\textquotedblright{} (Monier-Williams): \emph{a-} on \emph{nidrā}, sleep.
\end{cousinweb}

\begin{grammarlens}[title={What you've built}]
One more word for describing what you lived and felt.
\end{grammarlens}

\subsection*{Your turn}
\begin{itemize}
\raggedright
  \item \emph{Say it:} \emph{ônidrā}
  \item \emph{Say it:} \emph{ônidrā}, once more
  \item \emph{Say it:} \emph{chintāy bābā ônidrāy bhugchhen.} — Father is suffering from insomnia because of worry.
  \item \emph{Recall:} say the Bengali for to nurse, then the Bengali for sleep, then say \emph{ônidrā} again
\end{itemize}

\begin{tcolorbox}[breakable,colback=teal!4,colframe=teal!35!black,title={Before you move on}]
What does \bn{অনিদ্রা} mean? (\textquotedblleft{}Sleeplessness; insomnia\textquotedblright{}.) Say it once more.
\end{tcolorbox}
\section[\texorpdfstring{\bn{তন্দ্রা} (tôndrā)}{তন্দ্রা (tôndrā)}]{\bn{তন্দ্রা} (tôndrā) — drowsiness; a light doze}

\label{lesson:BN-C479-tondra}



\begin{quote}
Before the new one: say the Bengali for sleep, then the Bengali for insomnia.
\end{quote}

\subsection*{You'll want to know: \bn{তন্দ্রা}}
\textbf{\bn{তন্দ্রা}} — \emph{tôndrā} — \textquotedblleft{}drowsiness; a light doze\textquotedblright{}.

\begin{quote}
\bn{দুপুরে} \bn{খাওয়ার} \bn{পরে} \bn{আমার} \bn{তন্দ্রা} \bn{আসে।}
\end{quote}

\begin{quote}
\emph{dupure khāoyār pôre āmār tôndrā āse.}
\end{quote}

\begin{quote}
After lunch I get drowsy.
\end{quote}

\begin{cousinweb}
\textbf{\bn{তন্দ্রা}} (\emph{tôndrā}) is a tatsama, Sanskrit \emph{tandrā}, \textquotedblleft{}lassitude, exhaustion\textquotedblright{} (Monier-Williams); Bengali uses it for the drowsiness before sleep.
\end{cousinweb}

\begin{grammarlens}[title={What you've built}]
One more word for describing what you lived and felt.
\end{grammarlens}

\subsection*{Your turn}
\begin{itemize}
\raggedright
  \item \emph{Say it:} \emph{tôndrā}
  \item \emph{Say it:} \emph{tôndrā}, once more
  \item \emph{Say it:} \emph{dupure khāoyār pôre āmār tôndrā āse.} — After lunch I get drowsy.
  \item \emph{Recall:} say the Bengali for sleep, then the Bengali for insomnia, then say \emph{tôndrā} again
\end{itemize}

\begin{tcolorbox}[breakable,colback=teal!4,colframe=teal!35!black,title={Before you move on}]
What does \bn{তন্দ্রা} mean? (\textquotedblleft{}Drowsiness; a light doze\textquotedblright{}.) Say it once more.
\end{tcolorbox}
\section[\texorpdfstring{\bn{ঝিমানো} (jhimāno)}{ঝিমানো (jhimāno)}]{\bn{ঝিমানো} (jhimāno) — to doze; to nod off}

\label{lesson:BN-C479-jhimano}



\begin{quote}
Before the new one: say the Bengali for insomnia, then the Bengali for drowsiness.
\end{quote}

\subsection*{You'll want to know: \bn{ঝিমানো}}
\textbf{\bn{ঝিমানো}} — \emph{jhimāno} — \textquotedblleft{}to doze; to nod off\textquotedblright{}.

\begin{quote}
\bn{বাসে} \bn{বসে} \bn{দাদু} \bn{ঝিমাচ্ছিলেন।}
\end{quote}

\begin{quote}
\emph{bāse bôse dādu jhimāchchhilen.}
\end{quote}

\begin{quote}
Grandfather was nodding off on the bus.
\end{quote}

Past (he/she …): \bn{তিনি} \bn{চেয়ারে} \bn{ঝিমালেন}, \emph{tini cheyāre jhimālen}.

\begin{cousinweb}
\textbf{\bn{ঝিমানো}} (\emph{jhimāno}) is a Bengali verb for drowsy nodding whose origin is not settled.
\end{cousinweb}

\begin{grammarlens}[title={What you've built}]
One more word for describing what you lived and felt.
\end{grammarlens}

\subsection*{Your turn}
\begin{itemize}
\raggedright
  \item \emph{Say it:} \emph{jhimāno}
  \item \emph{Say it:} \emph{jhimāno}, once more
  \item \emph{Say it:} \emph{bāse bôse dādu jhimāchchhilen.} — Grandfather was nodding off on the bus.
  \item \emph{Recall:} say the Bengali for insomnia, then the Bengali for drowsiness, then say \emph{jhimāno} again
\end{itemize}

\begin{tcolorbox}[breakable,colback=teal!4,colframe=teal!35!black,title={Before you move on}]
What does \bn{ঝিমানো} mean? (\textquotedblleft{}To doze; to nod off\textquotedblright{}.) Say it once more.
\end{tcolorbox}
\section[\texorpdfstring{\bn{আড়মোড়া} \bn{ভাঙা} (āṛmoṛā bhāngā)}{আড়মোড়া ভাঙা (āṛmoṛā bhāngā)}]{\bn{আড়মোড়া} \bn{ভাঙা} (āṛmoṛā bhāngā) — to stretch on waking; to stretch one's limbs}

\label{lesson:BN-C479-armora-bhanga}



\begin{quote}
Before the new one: say the Bengali for drowsiness, then the Bengali for to nod off.
\end{quote}

\subsection*{You'll want to know: \bn{আড়মোড়া} \bn{ভাঙা}}
\textbf{\bn{আড়মোড়া} \bn{ভাঙা}} — \emph{āṛmoṛā bhāngā} — \textquotedblleft{}to stretch on waking; to stretch one's limbs\textquotedblright{}.

\begin{quote}
\bn{সকালে} \bn{জেগে} \bn{বিড়ালটা} \bn{আড়মোড়া} \bn{ভাঙল।}
\end{quote}

\begin{quote}
\emph{sôkāle jege biṛālṭā āṛmoṛā bhānglô.}
\end{quote}

\begin{quote}
Waking in the morning, the cat had a long stretch.
\end{quote}

Past (he/she …): \bn{তিনি} \bn{জেগে} \bn{আড়মোড়া} \bn{ভাঙলেন}, \emph{tini jege āṛmoṛā bhānglen}.

\begin{cousinweb}
\textbf{\bn{আড়মোড়া} \bn{ভাঙা}} (\emph{āṛmoṛā bhāngā}) is \emph{āṛmoṛā}, \textquotedblleft{}a stretch of the limbs\textquotedblright{}, whose origin is not settled, with \emph{bhāngā}, \textquotedblleft{}to break\textquotedblright{} (CDIAL 9353): to break out of the stiffness of sleep.
\end{cousinweb}

\begin{grammarlens}[title={What you've built}]
One more word for describing what you lived and felt.
\end{grammarlens}

\subsection*{Your turn}
\begin{itemize}
\raggedright
  \item \emph{Say it:} \emph{āṛmoṛā bhāngā}
  \item \emph{Say it:} \emph{āṛmoṛā bhāngā}, once more
  \item \emph{Say it:} \emph{sôkāle jege biṛālṭā āṛmoṛā bhānglô.} — Waking in the morning, the cat had a long stretch.
  \item \emph{Recall:} say the Bengali for drowsiness, then the Bengali for to nod off, then say \emph{āṛmoṛā bhāngā} again
\end{itemize}

\begin{tcolorbox}[breakable,colback=teal!4,colframe=teal!35!black,title={Before you move on}]
What does \bn{আড়মোড়া} \bn{ভাঙা} mean? (\textquotedblleft{}To stretch on waking; to stretch one's limbs\textquotedblright{}.) Say it once more.
\end{tcolorbox}
